\documentclass[10pt,a4paper]{amsart}
\usepackage{amsmath,amssymb,amsthm,graphicx}
\usepackage{verbatim}
\newenvironment{reference}{\comment}{\endcomment}
%\newenvironment{reference}{}{}
\newenvironment{slogan}{\comment}{\endcomment}
\newenvironment{history}{\comment}{\endcomment}

% For commutative diagrams we use Xy-pic
\usepackage[all]{xy}

% We use 2cell for 2-commutative diagrams.
\xyoption{2cell}
% Theorem environments.
%
\theoremstyle{plain}
\newtheorem{theorem}[subsection]{Theorem}
\newtheorem{proposition}[subsection]{Proposition}
\newtheorem{lemma}[subsection]{Lemma}

\theoremstyle{definition}
\newtheorem{definition}[subsection]{Definition}
\newtheorem{example}[subsection]{Example}
\newtheorem{exercise}[subsection]{Exercise}
\newtheorem{situation}[subsection]{Situation}

\theoremstyle{remark}
\newtheorem{remark}[subsection]{Remark}
\newtheorem{remarks}[subsection]{Remarks}

\numberwithin{equation}{subsection}

% Macros
%
\def\lim{\mathop{\mathrm{lim}}\nolimits}
\def\colim{\mathop{\mathrm{colim}}\nolimits}
\def\Spec{\mathop{\mathrm{Spec}}}
\def\Hom{\mathop{\mathrm{Hom}}\nolimits}
\def\Ext{\mathop{\mathrm{Ext}}\nolimits}
\def\SheafHom{\mathop{\mathcal{H}\!\mathit{om}}\nolimits}
\def\SheafExt{\mathop{\mathcal{E}\!\mathit{xt}}\nolimits}
\def\Sch{\mathit{Sch}}
\def\Mor{\mathop{\mathrm{Mor}}\nolimits}
\def\Ob{\mathop{\mathrm{Ob}}\nolimits}
\def\Sh{\mathop{\mathit{Sh}}\nolimits}
\def\NL{\mathop{N\!L}\nolimits}
\def\CH{\mathop{\mathrm{CH}}\nolimits}
\def\proetale{{pro\text{-}\acute{e}tale}}
\def\etale{{\acute{e}tale}}
\def\QCoh{\mathit{QCoh}}
\def\Ker{\mathop{\mathrm{Ker}}}
\def\Im{\mathop{\mathrm{Im}}}
\def\Coker{\mathop{\mathrm{Coker}}}
\def\Coim{\mathop{\mathrm{Coim}}}

% Boxtimes
%
\DeclareMathSymbol{\boxtimes}{\mathbin}{AMSa}{"02}

%
% Macros for moduli stacks/spaces
%
\def\QCohstack{\mathcal{QC}\!\mathit{oh}}
\def\Cohstack{\mathcal{C}\!\mathit{oh}}
\def\Spacesstack{\mathcal{S}\!\mathit{paces}}
\def\Quotfunctor{\mathrm{Quot}}
\def\Hilbfunctor{\mathrm{Hilb}}
\def\Curvesstack{\mathcal{C}\!\mathit{urves}}
\def\Polarizedstack{\mathcal{P}\!\mathit{olarized}}
\def\Complexesstack{\mathcal{C}\!\mathit{omplexes}}
% \Pic is the operator that assigns to X its picard group, usage \Pic(X)
% \Picardstack_{X/B} denotes the Picard stack of X over B
% \Picardfunctor_{X/B} denotes the Picard functor of X over B
\def\Pic{\mathop{\mathrm{Pic}}\nolimits}
\def\Picardstack{\mathcal{P}\!\mathit{ic}}
\def\Picardfunctor{\mathrm{Pic}}
\def\Deformationcategory{\mathcal{D}\!\mathit{ef}}
\usepackage{cleveref}
\crefname{originallemma}{original lemma}{original lemmas}
\Crefname{originallemma}{Original lemma}{Original lemmas}
\makeatletter
\newlabel{coherent-lemma-kill-by-twisting}{{02O1}{source}{}{sourcepointer-1}{}}
\newlabel{coherent-lemma-kill-by-twisting@cref}{{[originallemma][0][]02O1}{[1][0][]source}{}{}{}}
\makeatother
\crefname{originallemma}{original lemma}{original lemmas}
\Crefname{originallemma}{Original lemma}{Original lemmas}
\makeatletter
\newlabel{curves-lemma-tensor-omega-with-globally-generated-invertible}{{0B5H}{source}{}{sourcepointer-2}{}}
\newlabel{curves-lemma-tensor-omega-with-globally-generated-invertible@cref}{{[originallemma][0][]0B5H}{[1][0][]source}{}{}{}}
\makeatother
\crefname{originallemma}{original lemma}{original lemmas}
\Crefname{originallemma}{Original lemma}{Original lemmas}
\makeatletter
\newlabel{curves-lemma-vanishing-twist}{{0B62}{source}{}{sourcepointer-3}{}}
\newlabel{curves-lemma-vanishing-twist@cref}{{[originallemma][0][]0B62}{[1][0][]source}{}{}{}}
\makeatother
\crefname{originallemma}{original lemma}{original lemmas}
\Crefname{originallemma}{Original lemma}{Original lemmas}
\makeatletter
\newlabel{duality-lemma-sheaf-with-exact-support-effective-Cartier}{{0AA4}{source}{}{sourcepointer-4}{}}
\newlabel{duality-lemma-sheaf-with-exact-support-effective-Cartier@cref}{{[originallemma][0][]0AA4}{[1][0][]source}{}{}{}}
\makeatother
\crefname{originallemma}{original lemma}{original lemmas}
\Crefname{originallemma}{Original lemma}{Original lemmas}
\makeatletter
\newlabel{duality-lemma-twisted-inverse-image-closed}{{0A9X}{source}{}{sourcepointer-5}{}}
\newlabel{duality-lemma-twisted-inverse-image-closed@cref}{{[originallemma][0][]0A9X}{[1][0][]source}{}{}{}}
\makeatother
\crefname{originalreference}{original reference}{original references}
\Crefname{originalreference}{Original reference}{Original references}
\makeatletter
\newlabel{lemma-cohomology-blow-up-rational}{{0B63}{source}{}{sourcepointer-6}{}}
\newlabel{lemma-cohomology-blow-up-rational@cref}{{[originalreference][0][]0B63}{[1][0][]source}{}{}{}}
\makeatother
\crefname{originalreference}{original reference}{original references}
\Crefname{originalreference}{Original reference}{Original references}
\makeatletter
\newlabel{lemma-dualizing-rational}{{0BBU}{source}{}{sourcepointer-7}{}}
\newlabel{lemma-dualizing-rational@cref}{{[originalreference][0][]0BBU}{[1][0][]source}{}{}{}}
\makeatother
\crefname{originalreference}{original reference}{original references}
\Crefname{originalreference}{Original reference}{Original references}
\makeatletter
\newlabel{situation-rational}{{0B4W}{source}{}{sourcepointer-8}{}}
\newlabel{situation-rational@cref}{{[originalreference][0][]0B4W}{[1][0][]source}{}{}{}}
\makeatother
\title{Vanishing, global generation and surjectivity for the dualizing module of the blowup in the rational situation}
\author{The Stacks Project Authors}
\date{stacks-project:resolve.tex:local-snapshot-20261012; GNU FDL 1.2 or later}
\begin{document}
\maketitle
\noindent\textit{Editorial note.} All delivered mathematical text is the author's own, copied from the pinned Stacks Project chapter file resolve.tex; no statement or proof step was rewritten, reordered or removed.\par\medskip
\noindent\textit{Editorial note.} The counted proof is the author's complete argument: the identification of the restriction of the dualizing module to the exceptional divisor, the descending induction giving the vanishing, and the global generation argument yielding surjectivity.\par\medskip
\noindent\textit{Editorial note.} The label lemma-dualizing-blow-up-rational names the counted statement, and the proof environment immediately following it in the source is delivered verbatim as its proof.\par\medskip
\noindent\textit{Editorial note.} The counted statement is carried by the author's own lemma environment in the chapter source.\par\medskip
\noindent\textit{Editorial note.} Pre-open gate: python operations/agent-authoring/tools/precard\_gate.py --source resolve.tex --label lemma-dualizing-blow-up-rational --start 2640 --end 2685 --local-id 28103 returned RESULT=OK (L28103 free, occupancy FREE, public interval clear). Pre-assembly tag self-check: canonical tag resolved and not present in reference/public-items.jsonl.\par\medskip
\noindent\textit{Editorial note.} The statement and proof are copied from the pinned chapter file; the cited inputs are kept as external pointers with their canonical Stacks tags.\par\medskip
\noindent\textit{Allowed standard prerequisites (editorial).}\begin{itemize}
\item Cohomology of the blowup with values in the structure sheaf. Same chapter, the computation of the cohomology of the twists of the structure sheaf on the blowup of the maximal ideal of a regular local surface. Provides the vanishing input used to run the descending chain of cohomology groups.
\item Dualizing module of an effective Cartier divisor and its twist. Duality for Schemes, the lemmas expressing the dualizing module of an effective Cartier divisor and the twisted inverse image along a closed immersion. Identifies the restriction of the dualizing module of the blowup to the exceptional divisor with a twist of the dualizing module of the divisor.
\item Vanishing of the first cohomology of the dualizing sheaf on a curve, and global generation. Algebraic Curves, the lemma on the vanishing of the first cohomology of twists of the dualizing sheaf and the lemma making the dualizing sheaf tensored with a globally generated invertible sheaf globally generated. Supplies the surjectivity of the maps between the first cohomology groups and the global generation of the restriction to the exceptional divisor.
\end{itemize}\medskip
\ifdefined\hypertarget\hypertarget{sourcepointer-1}{}\fi
\noindent\textit{Source reference.} Original lemma 02O1: sources/stacks-chow/tags.txt, source lines 3500--3500. Prerequisite pointer invoked at proof line(s) 2673; it names a standard fact of a sibling chapter that the counted argument uses, and every step of the counted proof built on it is retained verbatim in the delivered body.\par\medskip
\ifdefined\hypertarget\hypertarget{sourcepointer-2}{}\fi
\noindent\textit{Source reference.} Original lemma 0B5H: sources/stacks-chow/tags.txt, source lines 14035--14035. Prerequisite pointer invoked at proof line(s) 2678; it names a standard fact of a sibling chapter that the counted argument uses, and every step of the counted proof built on it is retained verbatim in the delivered body.\par\medskip
\ifdefined\hypertarget\hypertarget{sourcepointer-3}{}\fi
\noindent\textit{Source reference.} Original lemma 0B62: sources/stacks-chow/tags.txt, source lines 14055--14055. Prerequisite pointer invoked at proof line(s) 2666; it names a standard fact of a sibling chapter that the counted argument uses, and every step of the counted proof built on it is retained verbatim in the delivered body.\par\medskip
\ifdefined\hypertarget\hypertarget{sourcepointer-4}{}\fi
\noindent\textit{Source reference.} Original lemma 0AA4: sources/stacks-chow/tags.txt, source lines 12972--12972. Prerequisite pointer invoked at proof line(s) 2660; it names a standard fact of a sibling chapter that the counted argument uses, and every step of the counted proof built on it is retained verbatim in the delivered body.\par\medskip
\ifdefined\hypertarget\hypertarget{sourcepointer-5}{}\fi
\noindent\textit{Source reference.} Original lemma 0A9X: sources/stacks-chow/tags.txt, source lines 12965--12965. Prerequisite pointer invoked at proof line(s) 2661; it names a standard fact of a sibling chapter that the counted argument uses, and every step of the counted proof built on it is retained verbatim in the delivered body.\par\medskip
\ifdefined\hypertarget\hypertarget{sourcepointer-6}{}\fi
\noindent\textit{Source reference.} Original reference 0B63: sources/stacks-resolve/pinned-original-source.tex, source lines 2581--2581. Prerequisite pointer invoked at proof line(s) 2656; it names a standard fact of the same chapter that the counted argument uses, and every step of the counted proof built on it is retained verbatim in the delivered body.\par\medskip
\ifdefined\hypertarget\hypertarget{sourcepointer-7}{}\fi
\noindent\textit{Source reference.} Original reference 0BBU: sources/stacks-resolve/pinned-original-source.tex, source lines 2618--2618. Prerequisite pointer invoked at proof line(s) 2651; it names a standard fact of the same chapter that the counted argument uses, and every step of the counted proof built on it is retained verbatim in the delivered body.\par\medskip
\ifdefined\hypertarget\hypertarget{sourcepointer-8}{}\fi
\noindent\textit{Source reference.} Original reference 0B4W: sources/stacks-resolve/pinned-original-source.tex, source lines 2455--2455. Prerequisite pointer invoked at proof line(s) 2642; it names a standard fact of the same chapter that the counted argument uses, and every step of the counted proof built on it is retained verbatim in the delivered body.\par\medskip
\begin{lemma}
\label{lemma-dualizing-blow-up-rational}
In Situation \ref{situation-rational} assume $A$ has a
dualizing complex $\omega_A^\bullet$ and is not regular.
Let $X$ be the blowup of $\Spec(A)$ in $\mathfrak m$ with
exceptional divisor $E \subset X$. Let $\omega_X$
be the dualizing module of $X$. Then
\begin{enumerate}
\item $\omega_E = \omega_X|_E \otimes \mathcal{O}_E(-1)$,
\item $H^1(X, \omega_X(n)) = 0$ for $n \geq 0$,
\item the map $f^*\omega_A \to \omega_X$ of
Lemma \ref{lemma-dualizing-rational} is surjective.
\end{enumerate}
\end{lemma}

\begin{proof}
We will use the results of Lemma \ref{lemma-cohomology-blow-up-rational}
without further mention. Observe that
$\omega_E = \omega_X|_E \otimes \mathcal{O}_E(-1)$
by Duality for Schemes, Lemmas
\ref{duality-lemma-sheaf-with-exact-support-effective-Cartier} and
\ref{duality-lemma-twisted-inverse-image-closed}. Thus
$\omega_X|_E = \omega_E(1)$. Consider the short exact sequences
$$
0 \to \omega_X(n + 1) \to \omega_X(n) \to i_*\omega_E(n + 1) \to 0
$$
By Algebraic Curves, Lemma \ref{curves-lemma-vanishing-twist}
we see that $H^1(E, \omega_E(n + 1)) = 0$ for $n \geq 0$.
Thus we see that the maps
$$
\ldots \to H^1(X, \omega_X(2)) \to H^1(X, \omega_X(1)) \to H^1(X, \omega_X)
$$
are surjective. Since $H^1(X, \omega_X(n))$ is zero for $n \gg 0$
(Cohomology of Schemes, Lemma \ref{coherent-lemma-kill-by-twisting})
we conclude that (2) holds.

\medskip\noindent
By Algebraic Curves, Lemma
\ref{curves-lemma-tensor-omega-with-globally-generated-invertible}
we see that $\omega_X|_E = \omega_E \otimes \mathcal{O}_E(1)$
is globally generated. Since we seen above that
$H^1(X, \omega_X(1)) = 0$ the map $H^0(X, \omega_X) \to H^0(E, \omega_X|_E)$
is surjective. We conclude that $\omega_X$ is globally generated
hence (3) holds because $\Gamma(X, \omega_X) = \omega_A$ is used
in Lemma \ref{lemma-dualizing-rational} to define the map.
\end{proof}

\end{document}
